\section*{5. Glossary EN $\leftrightarrow$ VI / Từ vựng song ngữ}
\scriptsize
\begin{multicols}{4}
\raggedright
\setlength{\parskip}{0pt}
\setlength{\parindent}{0pt}

\textbf{General CG / Chung}\\
\textbf{computer graphics} -- đồ họa máy tính\\
\textbf{rendering} -- kết xuất / dựng hình\\
\textbf{real-time rendering} -- kết xuất thời gian thực\\
\textbf{offline rendering} -- kết xuất ngoại tuyến\\
\textbf{ray tracing} -- dò tia\\
\textbf{ray casting} -- phóng tia\\
\textbf{path tracing} -- dò đường truyền tia\\
\textbf{rasterization} -- quét mành / điểm hóa\\
\textbf{pipeline} -- ống dẫn / quy trình xử lý\\
\textbf{graphics pipeline} -- ống dẫn đồ họa\\
\textbf{vertex} -- đỉnh\\
\textbf{edge} -- cạnh\\
\textbf{face} -- mặt\\
\textbf{polygon} -- đa giác\\
\textbf{primitive} -- nguyên thủy / khối cơ bản\\
\textbf{frame buffer} -- bộ đệm khung hình\\
\textbf{depth buffer / Z-buffer} -- bộ đệm độ sâu\\
\textbf{stencil buffer} -- bộ đệm khuôn\\
\textbf{shader} -- bộ đổ bóng / shader\\
\textbf{vertex shader} -- shader đỉnh\\
\textbf{fragment shader / pixel shader} -- shader điểm ảnh\\
\textbf{geometry shader} -- shader hình học\\
\textbf{compute shader} -- shader tính toán\\
\textbf{fragment} -- mảnh / phân đoạn\\
\textbf{pixel} -- điểm ảnh\\
\textbf{texel} -- điểm họa tiết\\
\textbf{voxel} -- điểm thể tích\\
\textbf{viewport} -- khung nhìn / cổng nhìn\\
\textbf{clipping} -- cắt / xén\\
\textbf{culling} -- loại bỏ\\
\textbf{back-face culling} -- loại bỏ mặt sau\\
\textbf{frustum culling} -- loại bỏ ngoài hình chóp cụt\\
\textbf{occlusion culling} -- loại bỏ bị che khuất\\
\textbf{projection} -- phép chiếu\\
\textbf{perspective projection} -- chiếu phối cảnh\\
\textbf{orthographic projection} -- chiếu trực giao / song song\\
\textbf{transformation} -- phép biến đổi\\
\textbf{rotation} -- phép quay\\
\textbf{translation} -- phép tịnh tiến\\
\textbf{scaling} -- phép co giãn / tỉ lệ\\
\textbf{shearing} -- phép trượt\\
\textbf{reflection} -- phép phản chiếu\\
\textbf{homogeneous coordinate} -- tọa độ thuần nhất\\
\textbf{world coordinate} -- tọa độ thế giới\\
\textbf{camera coordinate} -- tọa độ máy ảnh\\
\textbf{screen coordinate} -- tọa độ màn hình\\
\textbf{model-view matrix} -- ma trận mô hình-quan sát\\
\textbf{OpenGL} -- thư viện đồ họa OpenGL\\
\textbf{DirectX} -- thư viện DirectX\\
\textbf{Vulkan} -- API Vulkan\\
\textbf{GPU} -- bộ xử lý đồ họa\\
\textbf{CPU} -- bộ xử lý trung tâm\\
\textbf{anti-aliasing (AA)} -- khử răng cưa\\
\textbf{MSAA} -- khử răng cưa đa mẫu\\
\textbf{level of detail (LOD)} -- mức độ chi tiết\\
\textbf{texture} -- họa tiết / vân bề mặt\\
\textbf{texture mapping} -- ánh xạ họa tiết\\
\textbf{UV mapping} -- ánh xạ tọa độ UV\\
\textbf{mipmap} -- mipmap / đa mức họa tiết\\
\textbf{bump mapping} -- ánh xạ độ gồ ghề\\
\textbf{normal mapping} -- ánh xạ pháp tuyến\\
\textbf{displacement mapping} -- ánh xạ dịch chuyển\\
\textbf{normal vector} -- vec-tơ pháp tuyến\\
\textbf{tangent} -- tiếp tuyến\\
\textbf{binormal / bitangent} -- vec-tơ phụ tuyến\\
\textbf{lighting} -- chiếu sáng\\
\textbf{ambient light} -- ánh sáng môi trường\\
\textbf{diffuse light} -- ánh sáng khuếch tán\\
\textbf{specular light} -- ánh sáng phản chiếu\\
\textbf{point light} -- nguồn sáng điểm\\
\textbf{directional light} -- nguồn sáng định hướng\\
\textbf{spot light} -- nguồn sáng pha\\
\textbf{shading} -- tạo bóng / đổ bóng\\
\textbf{flat shading} -- đổ bóng phẳng\\
\textbf{Gouraud shading} -- đổ bóng Gouraud\\
\textbf{Phong shading} -- đổ bóng Phong\\
\textbf{PBR} -- kết xuất dựa trên vật lý\\
\textbf{material} -- chất liệu / vật liệu\\
\textbf{albedo} -- phản chiếu khuếch tán cơ bản\\
\textbf{roughness} -- độ nhám\\
\textbf{metallic} -- độ kim loại\\
\textbf{shadow mapping} -- ánh xạ bóng đổ\\
\textbf{global illumination} -- chiếu sáng toàn cục\\
\textbf{ambient occlusion (AO)} -- che khuất môi trường\\
\textbf{HDR} -- dải tương phản cao\\
\textbf{tone mapping} -- ánh xạ tông màu\\
\textbf{gamma correction} -- hiệu chỉnh gamma\\
\textbf{double buffering} -- đệm kép\\
\textbf{vsync} -- đồng bộ dọc\\
\textbf{frame rate (FPS)} -- tốc độ khung hình\\

\textbf{Geometric Modeling}\\
\textbf{geometric modeling} -- mô hình hóa hình học\\
\textbf{point cloud} -- đám mây điểm\\
\textbf{triangle mesh} -- lưới tam giác\\
\textbf{polygon mesh} -- lưới đa giác\\
\textbf{quad mesh} -- lưới tứ giác\\
\textbf{wireframe} -- khung dây\\
\textbf{manifold} -- đa tạp\\
\textbf{non-manifold} -- không đa tạp\\
\textbf{boundary} -- biên\\
\textbf{boundary edge} -- cạnh biên\\
\textbf{hole} -- lỗ thủng\\
\textbf{hole filling} -- vá lỗ / lấp lỗ\\
\textbf{watertight} -- kín nước / liền khối\\
\textbf{closed surface} -- mặt kín\\
\textbf{open surface} -- mặt hở\\
\textbf{orientable} -- có hướng\\
\textbf{genus} -- chi (số lỗ)\\
\textbf{Euler characteristic} -- đặc trưng Euler\\
\textbf{topology} -- tô-pô / hình học liên kết\\
\textbf{geometry} -- hình học (tọa độ)\\
\textbf{half-edge} -- nửa cạnh\\
\textbf{winged-edge} -- cạnh có cánh\\
\textbf{adjacency} -- kề / lân cận\\
\textbf{normal} -- pháp tuyến\\
\textbf{vertex normal} -- pháp tuyến đỉnh\\
\textbf{face normal} -- pháp tuyến mặt\\
\textbf{surface} -- bề mặt\\
\textbf{curve} -- đường cong\\
\textbf{simplification} -- đơn giản hóa\\
\textbf{decimation} -- giảm lưới / giảm mật độ\\
\textbf{mesh simplification} -- đơn giản hóa lưới\\
\textbf{remeshing} -- tạo lại lưới\\
\textbf{triangulation} -- tam giác hóa\\
\textbf{Delaunay triangulation} -- tam giác hóa Delaunay\\
\textbf{constrained Delaunay} -- Delaunay có ràng buộc\\
\textbf{Voronoi diagram} -- biểu đồ Voronoi\\
\textbf{circumcircle} -- đường tròn ngoại tiếp\\
\textbf{in-circle predicate} -- vị từ trong đường tròn\\
\textbf{empty circle property} -- tính chất đường tròn rỗng\\
\textbf{convex hull} -- bao lồi\\
\textbf{flip operation / edge flip} -- phép lật cạnh\\
\textbf{fairing} -- làm mượt / chỉnh trơn\\
\textbf{smoothing} -- làm trơn\\
\textbf{denoising} -- khử nhiễu\\
\textbf{noise removal} -- loại bỏ nhiễu\\
\textbf{outlier} -- điểm ngoại lai\\
\textbf{sliver triangle} -- tam giác mỏng / tam giác kim\\
\textbf{regular sampling} -- lấy mẫu đều\\
\textbf{irregular sampling} -- lấy mẫu không đều\\
\textbf{alpha shape} -- hình dạng alpha\\
\textbf{ball pivoting (BPA)} -- xoay bóng\\
\textbf{Poisson reconstruction} -- tái tạo Poisson\\
\textbf{Marching Cubes} -- thuật toán Marching Cubes\\
\textbf{isosurface} -- mặt đẳng trị\\
\textbf{isovalue} -- giá trị đẳng trị\\
\textbf{scalar field} -- trường vô hướng\\
\textbf{signed distance function (SDF)} -- hàm khoảng cách có dấu\\
\textbf{implicit surface} -- mặt ẩn\\
\textbf{explicit surface} -- mặt tường minh\\
\textbf{parametric surface} -- mặt tham số\\
\textbf{NURBS} -- B-spline hữu tỷ không đều\\
\textbf{B-spline} -- đường spline cơ sở B\\
\textbf{Bezier} -- đường Bezier\\
\textbf{control point} -- điểm điều khiển\\
\textbf{control polygon} -- đa giác điều khiển\\
\textbf{control net} -- lưới điều khiển\\
\textbf{knot} -- nút\\
\textbf{knot vector} -- vec-tơ nút\\
\textbf{basis function} -- hàm cơ sở\\
\textbf{tensor product} -- tích tensor\\
\textbf{subdivision surface} -- mặt chia nhỏ\\
\textbf{subdivision} -- chia nhỏ / mịn hóa\\
\textbf{Catmull-Clark} -- Catmull-Clark\\
\textbf{Loop subdivision} -- chia nhỏ Loop\\
\textbf{refinement} -- mịn hóa / tinh chỉnh\\
\textbf{edge collapse} -- gộp cạnh\\
\textbf{vertex split} -- tách đỉnh\\
\textbf{quadric error metric (QEM)} -- độ đo sai số bậc hai\\
\textbf{Laplacian} -- toán tử Laplace\\
\textbf{Laplacian smoothing} -- làm trơn Laplace\\
\textbf{cotangent weight} -- trọng số cotang\\
\textbf{Taubin smoothing} -- làm trơn Taubin\\
\textbf{bilateral filter} -- bộ lọc song phương\\
\textbf{Gaussian filter} -- bộ lọc Gauss\\
\textbf{mean curvature} -- độ cong trung bình\\
\textbf{Gaussian curvature} -- độ cong Gauss\\
\textbf{principal curvature} -- độ cong chính\\
\textbf{moving least squares (MLS)} -- bình phương tối thiểu di động\\
\textbf{WLOP} -- vận hành định vị địa phương có trọng số\\
\textbf{PCA} -- phân tích thành phần chính\\
\textbf{kNN} -- k láng giềng gần nhất\\
\textbf{octree} -- cây bát phân\\
\textbf{KD-tree} -- cây KD\\
\textbf{BVH} -- phân cấp bao đóng\\
\textbf{bounding box} -- hộp bao\\
\textbf{AABB} -- hộp bao trục căn chỉnh\\
\textbf{OBB} -- hộp bao có hướng\\
\textbf{registration} -- căn chỉnh / khớp lưới\\
\textbf{ICP (Iterative Closest Point)} -- thuật toán điểm gần nhất lặp\\
\textbf{Hausdorff distance} -- khoảng cách Hausdorff\\
\textbf{advancing front} -- mặt tiến\\
\textbf{Greedy projection} -- chiếu tham lam\\

\textbf{Medical Imaging}\\
\textbf{medical image} -- ảnh y tế\\
\textbf{MRI} -- chụp cộng hưởng từ\\
\textbf{CT scan} -- chụp cắt lớp vi tính\\
\textbf{ultrasound (US)} -- siêu âm\\
\textbf{PET} -- chụp cắt lớp phát xạ positron\\
\textbf{X-ray} -- tia X / chụp X-quang\\
\textbf{SPECT} -- chụp cắt lớp đơn photon\\
\textbf{OCT} -- chụp cắt lớp quang học\\
\textbf{fMRI} -- cộng hưởng từ chức năng\\
\textbf{DICOM} -- chuẩn ảnh y tế DICOM\\
\textbf{NIfTI} -- định dạng NIfTI\\
\textbf{modality} -- phương thức / phương pháp ảnh\\
\textbf{slice} -- lát cắt\\
\textbf{volume} -- khối thể tích\\
\textbf{volume rendering} -- kết xuất khối\\
\textbf{multi-planar reformatting (MPR)} -- tái dựng đa mặt phẳng\\
\textbf{k-space} -- không gian k\\
\textbf{sinogram} -- ảnh sin / sinogram\\
\textbf{Radon transform} -- biến đổi Radon\\
\textbf{Filtered Back Projection (FBP)} -- chiếu ngược có lọc\\
\textbf{Hounsfield Unit (HU)} -- đơn vị Hounsfield\\
\textbf{T1-weighted} -- ảnh trọng số T1\\
\textbf{T2-weighted} -- ảnh trọng số T2\\
\textbf{FLAIR} -- ảnh FLAIR\\
\textbf{DWI} -- ảnh khuếch tán\\
\textbf{ADC} -- hệ số khuếch tán biểu kiến\\
\textbf{contrast agent} -- thuốc cản quang / chất tương phản\\
\textbf{gadolinium} -- gadolini\\
\textbf{iodine contrast} -- cản quang iốt\\
\textbf{segmentation} -- phân đoạn\\
\textbf{semantic segmentation} -- phân đoạn ngữ nghĩa\\
\textbf{instance segmentation} -- phân đoạn đối tượng\\
\textbf{registration} -- khớp ảnh / căn chỉnh ảnh\\
\textbf{rigid registration} -- khớp ảnh cứng\\
\textbf{non-rigid registration} -- khớp ảnh phi cứng\\
\textbf{classification} -- phân loại\\
\textbf{detection} -- phát hiện\\
\textbf{brain tumor} -- u não\\
\textbf{glioma} -- u thần kinh đệm\\
\textbf{meningioma} -- u màng não\\
\textbf{BraTS} -- thử thách phân đoạn u não BraTS\\
\textbf{lesion} -- tổn thương\\
\textbf{edema} -- phù nề\\
\textbf{hemorrhage} -- xuất huyết\\
\textbf{necrosis} -- hoại tử\\
\textbf{enhancing tumor} -- u ngấm thuốc\\
\textbf{whole tumor} -- toàn bộ u\\
\textbf{tumor core} -- lõi u\\
\textbf{skull stripping} -- loại bỏ hộp sọ\\
\textbf{bias field} -- nhiễu trường thiên lệch\\
\textbf{N4 bias correction} -- hiệu chỉnh thiên lệch N4\\
\textbf{intensity normalization} -- chuẩn hóa cường độ\\
\textbf{histogram matching} -- khớp biểu đồ\\
\textbf{Otsu thresholding} -- ngưỡng Otsu\\
\textbf{thresholding} -- phân ngưỡng\\
\textbf{watershed} -- phân vùng đường phân thủy\\
\textbf{region growing} -- phát triển vùng\\
\textbf{snake / active contour} -- đường viền chủ động\\
\textbf{level set} -- tập mức\\
\textbf{morphological operation} -- phép toán hình thái\\
\textbf{erosion} -- co\\
\textbf{dilation} -- giãn\\
\textbf{opening} -- mở\\
\textbf{closing} -- đóng\\
\textbf{U-Net} -- mạng U-Net\\
\textbf{V-Net} -- mạng V-Net\\
\textbf{nnU-Net} -- mạng nnU-Net (tự cấu hình)\\
\textbf{encoder-decoder} -- bộ mã hóa-giải mã\\
\textbf{skip connection} -- kết nối tắt\\
\textbf{Dice coefficient} -- hệ số Dice\\
\textbf{Jaccard / IoU} -- chỉ số Jaccard / giao trên hợp\\
\textbf{Hausdorff distance} -- khoảng cách Hausdorff\\
\textbf{ASSD} -- khoảng cách bề mặt đối xứng trung bình\\
\textbf{sensitivity / recall} -- độ nhạy\\
\textbf{specificity} -- độ đặc hiệu\\
\textbf{precision} -- độ chính xác (precision)\\
\textbf{accuracy} -- độ chính xác tổng\\
\textbf{F1 score} -- điểm F1\\
\textbf{ROC curve} -- đường ROC\\
\textbf{AUC} -- diện tích dưới đường cong\\
\textbf{ground truth} -- nhãn thật / nhãn chuẩn\\
\textbf{label} -- nhãn\\
\textbf{annotation} -- chú thích\\
\textbf{augmentation} -- tăng cường dữ liệu\\
\textbf{transfer learning} -- học chuyển giao\\
\textbf{fine-tuning} -- tinh chỉnh\\
\textbf{pretrained model} -- mô hình đã huấn luyện sẵn\\
\textbf{GAN} -- mạng đối kháng sinh\\
\textbf{generator} -- bộ sinh\\
\textbf{discriminator} -- bộ phân biệt\\
\textbf{cGAN} -- GAN có điều kiện\\
\textbf{diagnosis} -- chẩn đoán\\
\textbf{prognosis} -- tiên lượng\\
\textbf{CDSS} -- hệ thống hỗ trợ quyết định lâm sàng\\
\textbf{XAI} -- AI giải thích được\\
\textbf{Grad-CAM} -- bản đồ kích hoạt theo gradient\\
\textbf{saliency map} -- bản đồ nổi bật\\
\textbf{attention map} -- bản đồ chú ý\\

\textbf{VR / AR / XR}\\
\textbf{virtual reality (VR)} -- thực tế ảo\\
\textbf{augmented reality (AR)} -- thực tế tăng cường\\
\textbf{mixed reality (MR)} -- thực tế hỗn hợp\\
\textbf{extended reality (XR)} -- thực tế mở rộng\\
\textbf{HMD / head-mounted display} -- kính/mũ đeo đầu\\
\textbf{immersion} -- sự đắm chìm / nhập vai\\
\textbf{presence} -- cảm giác hiện diện\\
\textbf{telepresence} -- hiện diện từ xa\\
\textbf{stereoscopic} -- lập thể\\
\textbf{stereo rendering} -- kết xuất lập thể\\
\textbf{parallax} -- thị sai\\
\textbf{depth perception} -- nhận thức độ sâu\\
\textbf{binocular disparity} -- chênh lệch hai mắt\\
\textbf{accommodation} -- điều tiết\\
\textbf{vergence} -- hội tụ mắt\\
\textbf{vergence-accommodation conflict} -- mâu thuẫn hội tụ-điều tiết\\
\textbf{FOV (field of view)} -- góc nhìn\\
\textbf{IPD (interpupillary distance)} -- khoảng cách đồng tử\\
\textbf{refresh rate} -- tốc độ làm mới\\
\textbf{latency} -- độ trễ\\
\textbf{motion-to-photon latency} -- độ trễ chuyển động đến photon\\
\textbf{tracking} -- theo dõi / bám\\
\textbf{6DoF} -- 6 bậc tự do\\
\textbf{3DoF} -- 3 bậc tự do\\
\textbf{inside-out tracking} -- theo dõi từ trong ra\\
\textbf{outside-in tracking} -- theo dõi từ ngoài vào\\
\textbf{lighthouse} -- trạm phát Lighthouse\\
\textbf{IMU} -- đơn vị đo quán tính\\
\textbf{accelerometer} -- gia tốc kế\\
\textbf{gyroscope} -- con quay hồi chuyển\\
\textbf{magnetometer} -- từ kế\\
\textbf{SLAM} -- định vị và lập bản đồ đồng thời\\
\textbf{visual SLAM} -- SLAM thị giác\\
\textbf{marker-based AR} -- AR dựa trên dấu mốc\\
\textbf{markerless AR} -- AR không dấu mốc\\
\textbf{fiducial marker} -- dấu mốc nhận dạng\\
\textbf{ArUco / QR marker} -- dấu ArUco / QR\\
\textbf{anchor} -- điểm neo\\
\textbf{plane detection} -- phát hiện mặt phẳng\\
\textbf{feature point} -- điểm đặc trưng\\
\textbf{occlusion} -- che khuất\\
\textbf{pass-through} -- nhìn xuyên (camera)\\
\textbf{optical see-through} -- nhìn xuyên quang học\\
\textbf{video see-through} -- nhìn xuyên video\\
\textbf{waveguide} -- ống dẫn sóng\\
\textbf{OLED} -- màn hình OLED\\
\textbf{micro-OLED} -- micro-OLED\\
\textbf{LCD} -- màn hình LCD\\
\textbf{foveated rendering} -- kết xuất theo điểm nhìn\\
\textbf{eye tracking} -- theo dõi mắt\\
\textbf{lens distortion} -- méo thấu kính\\
\textbf{barrel distortion} -- méo dạng thùng\\
\textbf{pincushion distortion} -- méo dạng gối\\
\textbf{chromatic aberration} -- quang sai sắc\\
\textbf{motion sickness} -- say tàu xe ảo\\
\textbf{cybersickness} -- say không gian mạng\\
\textbf{vection} -- ảo giác chuyển động\\
\textbf{controller} -- tay điều khiển\\
\textbf{haptic feedback} -- phản hồi xúc giác\\
\textbf{force feedback} -- phản hồi lực\\
\textbf{ray casting (VR)} -- phóng tia chọn vật\\
\textbf{teleportation} -- dịch chuyển tức thời\\
\textbf{snap turn} -- xoay nhanh\\
\textbf{smooth locomotion} -- di chuyển trơn\\
\textbf{room scale} -- quy mô phòng\\
\textbf{guardian / chaperone} -- vùng an toàn\\
\textbf{Unity} -- engine Unity\\
\textbf{Unreal Engine} -- engine Unreal\\
\textbf{ARKit} -- ARKit (Apple)\\
\textbf{ARCore} -- ARCore (Google)\\
\textbf{Vuforia} -- Vuforia\\
\textbf{OpenXR} -- chuẩn OpenXR\\
\textbf{MRTK} -- bộ công cụ MR (Microsoft)\\
\textbf{WebXR} -- WebXR (XR trên web)\\
\textbf{eMuseum} -- bảo tàng điện tử\\
\textbf{eTourism} -- du lịch điện tử\\
\textbf{digital heritage} -- di sản số\\
\textbf{digital twin} -- bản sao số\\
\textbf{avatar} -- nhân vật đại diện\\
\textbf{metaverse} -- vũ trụ ảo\\

\textbf{Bezier \& Splines}\\
\textbf{Bezier curve} -- đường cong Bezier\\
\textbf{Bezier surface} -- mặt Bezier\\
\textbf{Bernstein polynomial} -- đa thức Bernstein\\
\textbf{Bernstein basis} -- cơ sở Bernstein\\
\textbf{control point} -- điểm điều khiển\\
\textbf{control polygon} -- đa giác điều khiển\\
\textbf{control net} -- lưới điều khiển\\
\textbf{parameter $t$} -- tham số $t$\\
\textbf{parameterization} -- tham số hóa\\
\textbf{interpolation} -- nội suy\\
\textbf{approximation} -- xấp xỉ\\
\textbf{De Casteljau algorithm} -- thuật toán De Casteljau\\
\textbf{recursion} -- đệ quy\\
\textbf{subdivision} -- chia đôi / chia nhỏ\\
\textbf{degree} -- bậc\\
\textbf{degree elevation} -- nâng bậc\\
\textbf{degree reduction} -- hạ bậc\\
\textbf{derivative} -- đạo hàm\\
\textbf{hodograph} -- đồ thị đạo hàm / hodograph\\
\textbf{tangent vector} -- vec-tơ tiếp tuyến\\
\textbf{normal vector} -- vec-tơ pháp tuyến\\
\textbf{curvature} -- độ cong\\
\textbf{torsion} -- độ xoắn\\
\textbf{convex hull property} -- tính chất bao lồi\\
\textbf{partition of unity} -- phân hoạch đơn vị\\
\textbf{affine invariance} -- bất biến affine\\
\textbf{variation diminishing} -- giảm biến thiên\\
\textbf{symmetry} -- đối xứng\\
\textbf{endpoint interpolation} -- nội suy điểm đầu cuối\\
\textbf{isoparametric curve} -- đường đẳng tham số\\
\textbf{tensor product surface} -- mặt tích tensor\\
\textbf{bicubic} -- bậc ba kép / song bậc ba\\
\textbf{biquadratic} -- bậc hai kép\\
\textbf{rational Bezier} -- Bezier hữu tỷ\\
\textbf{B-spline} -- B-spline\\
\textbf{NURBS} -- NURBS\\
\textbf{knot vector} -- vec-tơ nút\\
\textbf{uniform knot} -- nút đều\\
\textbf{non-uniform knot} -- nút không đều\\
\textbf{Hermite curve} -- đường Hermite\\
\textbf{Catmull-Rom spline} -- spline Catmull-Rom\\
\textbf{cubic spline} -- spline bậc ba\\
\textbf{spline} -- spline / đường ghép trơn\\
\textbf{continuity} -- tính liên tục\\
\textbf{C0 continuity} -- liên tục C0 (cùng vị trí)\\
\textbf{C1 continuity} -- liên tục đạo hàm bậc 1\\
\textbf{C2 continuity} -- liên tục đạo hàm bậc 2\\
\textbf{G1 continuity} -- liên tục hình học G1\\
\textbf{G2 continuity} -- liên tục hình học G2\\
\textbf{smoothness} -- độ trơn\\
\textbf{blending function} -- hàm pha trộn\\
\textbf{weight} -- trọng số\\
\textbf{rational basis} -- cơ sở hữu tỷ\\
\textbf{Coons patch} -- mảnh Coons\\
\textbf{Gregory patch} -- mảnh Gregory\\
\textbf{control mesh} -- lưới điều khiển\\
\textbf{patch} -- mảnh\\

\textbf{Math / Toán}\\
\textbf{matrix} -- ma trận\\
\textbf{vector} -- vec-tơ\\
\textbf{scalar} -- vô hướng\\
\textbf{determinant} -- định thức\\
\textbf{eigenvalue} -- giá trị riêng\\
\textbf{eigenvector} -- vec-tơ riêng\\
\textbf{inverse matrix} -- ma trận nghịch đảo\\
\textbf{transpose} -- chuyển vị\\
\textbf{identity matrix} -- ma trận đơn vị\\
\textbf{orthogonal matrix} -- ma trận trực giao\\
\textbf{symmetric matrix} -- ma trận đối xứng\\
\textbf{positive definite} -- xác định dương\\
\textbf{rank} -- hạng\\
\textbf{trace} -- vết\\
\textbf{dot product / inner product} -- tích vô hướng\\
\textbf{cross product} -- tích có hướng\\
\textbf{outer product} -- tích ngoài\\
\textbf{norm} -- chuẩn\\
\textbf{$L_2$ norm} -- chuẩn $L_2$ / chuẩn Euclid\\
\textbf{gradient} -- gradient / vec-tơ độ dốc\\
\textbf{divergence} -- tán xạ / div\\
\textbf{Laplacian} -- toán tử Laplace\\
\textbf{Hessian} -- ma trận Hess\\
\textbf{Jacobian} -- ma trận Jacobi\\
\textbf{partial derivative} -- đạo hàm riêng\\
\textbf{integral} -- tích phân\\
\textbf{Fourier transform} -- biến đổi Fourier\\
\textbf{FFT} -- biến đổi Fourier nhanh\\
\textbf{convolution} -- tích chập\\
\textbf{kernel} -- nhân (toán)\\
\textbf{basis} -- cơ sở\\
\textbf{span} -- bao tuyến tính\\
\textbf{linear combination} -- tổ hợp tuyến tính\\
\textbf{projection} -- phép chiếu\\
\textbf{orthogonal} -- trực giao\\
\textbf{orthonormal} -- trực chuẩn\\
\textbf{plane} -- mặt phẳng\\
\textbf{line} -- đường thẳng\\
\textbf{ray} -- tia\\
\textbf{intersection} -- giao\\
\textbf{distance} -- khoảng cách\\
\textbf{angle} -- góc\\
\textbf{area} -- diện tích\\
\textbf{volume} -- thể tích\\
\textbf{polynomial} -- đa thức\\
\textbf{recursion} -- đệ quy\\
\textbf{iteration} -- lặp\\
\textbf{convergence} -- hội tụ\\
\textbf{optimization} -- tối ưu hóa\\
\textbf{minimization} -- cực tiểu hóa\\
\textbf{maximization} -- cực đại hóa\\
\textbf{gradient descent} -- giảm theo gradient\\
\textbf{stochastic gradient descent} -- giảm gradient ngẫu nhiên\\
\textbf{regression} -- hồi quy\\
\textbf{regularization} -- chính quy hóa\\
\textbf{least squares} -- bình phương tối thiểu\\
\textbf{SVD} -- phân tích giá trị suy biến\\
\textbf{quaternion} -- số quaternion\\
\textbf{Euler angle} -- góc Euler\\
\textbf{barycentric coordinate} -- tọa độ trọng tâm\\

\textbf{Algorithms / Performance}\\
\textbf{algorithm} -- thuật toán\\
\textbf{complexity} -- độ phức tạp\\
\textbf{time complexity} -- độ phức tạp thời gian\\
\textbf{space complexity} -- độ phức tạp không gian\\
\textbf{big-O notation} -- ký hiệu O lớn\\
\textbf{recursion} -- đệ quy\\
\textbf{dynamic programming} -- quy hoạch động\\
\textbf{greedy algorithm} -- thuật toán tham lam\\
\textbf{divide and conquer} -- chia để trị\\
\textbf{backtracking} -- quay lui\\
\textbf{heap} -- đống / heap\\
\textbf{priority queue} -- hàng đợi ưu tiên\\
\textbf{stack} -- ngăn xếp\\
\textbf{queue} -- hàng đợi\\
\textbf{hash table} -- bảng băm\\
\textbf{BFS} -- duyệt theo chiều rộng\\
\textbf{DFS} -- duyệt theo chiều sâu\\
\textbf{Dijkstra} -- thuật toán Dijkstra\\
\textbf{A* search} -- tìm kiếm A sao\\
\textbf{sort} -- sắp xếp\\
\textbf{search} -- tìm kiếm\\
\textbf{binary search} -- tìm kiếm nhị phân\\
\textbf{graph} -- đồ thị\\
\textbf{tree} -- cây\\
\textbf{octree} -- cây bát phân\\
\textbf{quadtree} -- cây tứ phân\\
\textbf{KD-tree} -- cây KD\\
\textbf{R-tree} -- cây R\\
\textbf{BVH} -- phân cấp khối bao\\
\textbf{spatial index} -- chỉ mục không gian\\
\textbf{nearest neighbor} -- láng giềng gần nhất\\
\textbf{parallel computing} -- tính toán song song\\
\textbf{concurrency} -- đồng thời\\
\textbf{thread} -- luồng\\
\textbf{GPU computing} -- tính toán trên GPU\\
\textbf{CUDA} -- nền tảng CUDA\\
\textbf{OpenCL} -- nền tảng OpenCL\\
\textbf{SIMD} -- một lệnh nhiều dữ liệu\\
\textbf{cache} -- bộ nhớ đệm\\
\textbf{memory bandwidth} -- băng thông bộ nhớ\\
\textbf{benchmark} -- đánh giá hiệu năng\\
\textbf{profiling} -- đo lường hiệu năng\\
\textbf{optimization} -- tối ưu hóa\\
\textbf{lookup table (LUT)} -- bảng tra cứu\\
\textbf{precomputation} -- tính trước\\
\textbf{lazy evaluation} -- ước lượng trễ\\
\textbf{streaming} -- truyền dòng\\
\textbf{out-of-core} -- ngoài bộ nhớ chính\\

\textbf{Deep Learning (bonus)}\\
\textbf{neural network} -- mạng nơ-ron\\
\textbf{deep learning} -- học sâu\\
\textbf{CNN} -- mạng nơ-ron tích chập\\
\textbf{RNN} -- mạng nơ-ron hồi quy\\
\textbf{transformer} -- mạng transformer\\
\textbf{attention} -- cơ chế chú ý\\
\textbf{self-attention} -- tự chú ý\\
\textbf{layer} -- tầng / lớp\\
\textbf{convolution layer} -- tầng tích chập\\
\textbf{pooling} -- gộp / lấy mẫu xuống\\
\textbf{max pooling} -- gộp lớn nhất\\
\textbf{batch normalization} -- chuẩn hóa lô\\
\textbf{dropout} -- bỏ ngẫu nhiên\\
\textbf{activation function} -- hàm kích hoạt\\
\textbf{ReLU} -- hàm ReLU\\
\textbf{sigmoid} -- hàm sigmoid\\
\textbf{softmax} -- hàm softmax\\
\textbf{loss function} -- hàm mất mát\\
\textbf{cross-entropy} -- entropy chéo\\
\textbf{Dice loss} -- mất mát Dice\\
\textbf{backpropagation} -- lan truyền ngược\\
\textbf{epoch} -- chu kỳ huấn luyện\\
\textbf{batch} -- lô\\
\textbf{learning rate} -- tốc độ học\\
\textbf{overfitting} -- quá khớp\\
\textbf{underfitting} -- thiếu khớp\\
\textbf{validation set} -- tập kiểm định\\
\textbf{NeRF} -- trường bức xạ thần kinh\\
\textbf{diffusion model} -- mô hình khuếch tán\\

\end{multicols}
